\scanpage{386}\mspage{125}

en avoir. \deleted{\unsure{C'est}} Mais pas pour les $C_\infty(a,b)$ eux-mêmes ?

C'est un peu vexant qu'avec tout ce travail, il semblerait qu'on n'ait pas même été \unsure{foutus|foutu} de définir \inserted{l'analogue dans $\mathrm{Cat}$} \unsure{des classiques espaces|du classique espace} \unsure{des|de} lacets $\Omega(X,a)$ comme un groupe \uline{dans Hot}. (Que fait Thomason ?) Mais voilà une idée \unsure{pendable|pensable|possible}. Pour définir \uline{quand même}

\[
C_\infty(a,b)\times C_\infty(b,c)\unsure{-|\longrightarrow} C_\infty(a,c)
\]

dans $\mathrm{Hot}$, on va décrire les deux facteurs $C_\infty(a,b)$ et $C_\infty(b,c)$ par des systèmes projectifs différents $\underline t,\underline t'$. On \deleted{\unsure{prendra disons}} \inserted{peut l'expliciter} en partant d'un $\tau_0$ de longueur \unsure{utile} $\geq2$, et en prenant dans les deux cas

\[
t_n=t'_n=\tau_0^n,
\]

mais des morphismes de transition différents

\[
\begin{aligned}
(t_{n+p}\longrightarrow t_n)&=\pi_p*\mathrm{id}_{\tau_0^n}\\
(t'_{n+p}-t'_n)&=\mathrm{id}_{\tau_0^n}*\pi_p
\end{aligned}
\]

où $\pi_p:\tau_0^p-1$ est la \deleted{\unsure{can}} dégénérescence totale.

On \inserted{\unsure{donc}} définira $C_\infty(a,c)$ en termes de $\underline t''\inserted{=\underline t*\underline t'}$

\[
\left\{\begin{aligned}
t''_n&=\tau_0^{2n}\\
(t''_{n+p}-t''_n)&=\pi_p*\mathrm{id}_{\tau_0^{2n}}*\pi_{\unsure{p}}.
\end{aligned}\right.
\]

Cette fois, les \deleted{morphismes} \inserted{foncteurs} de composition

\[
\deleted{\unsure{H}}C_n(\deleted{\unsure{t_n}}\underline t;a,b)\times C_n(\underline t';\unsure{b,c})\longrightarrow C_n(\underline t'';a,c)
\]

sont compatibles avec les foncteurs transition pour $n$ variable, et on trouve

% archive sheet number: 385
